% Artykuł własny: LNCS, po polsku, IMRaD, max 6 stron bez bibliografii.
% Kompilacja: latexmk -pdf main.tex
% Tabele w tables/ generuje src/summarize.py -- nie edytować ręcznie.
% \TODO{...} oznacza miejsca do uzupełnienia przez grupę.
\documentclass[runningheads]{llncs}
\usepackage[T1]{fontenc}
\usepackage[utf8]{inputenc}
\usepackage{polski}
\usepackage{lmodern}
\usepackage{amsmath}
\usepackage{booktabs}
\usepackage{graphicx}
\usepackage{xcolor}
\usepackage{hyperref}
\newcommand{\TODO}[1]{\textcolor{red}{[TODO: #1]}}
\renewcommand{\abstractname}{Streszczenie}
\renewcommand{\keywordname}{\textbf{Słowa kluczowe:}}
\renewcommand{\tablename}{Tabela}
\renewcommand{\figurename}{Rys.}
\renewcommand{\refname}{Bibliografia}

\begin{document}

\title{Ocena silników OCR i modeli wizyjno-językowych w~rozpoznawaniu tekstu polskiego na danych syntetycznych}
\titlerunning{OCR i VLM dla tekstu polskiego}
\author{\TODO{Autor 1} \and \TODO{Autor 2} \and \TODO{Autor 3}}
\authorrunning{\TODO{Autorzy}}
\institute{\TODO{Uczelnia, wydział}\\\email{\TODO{e-mail}}}

\maketitle

\begin{abstract}
Systemy OCR i~modele wizyjno-językowe (VLM) trenowane są głównie na danych angielskich, przez co
gorzej radzą sobie ze znakami diakrytycznymi innych języków. Przenosimy protokół badania
Yılmaza i~in.~\cite{yilmaz2026} z~języka tureckiego na polski: generujemy syntetyczny zbiór obrazów
z~\emph{Lalki} B.~Prusa (tekst drukowany i~odręczny), kontrolując obecność polskich znaków, długość słowa,
kontekst zdania i~pięć rodzajów zniekształceń. Porównujemy \TODO{liczba} modeli: Tesseract i~EasyOCR
z~modelami polskimi, PaddleOCR, TrOCR oraz Qwen2.5-VL. \TODO{Główny wynik w~jednym zdaniu.} Pokazujemy
też, że słownik znaków PaddleOCR dla języków łacińskich nie zawiera liter ń, ś, ź, ż, co wyklucza
ich poprawne rozpoznanie. Kod i~dane udostępniamy publicznie.
\keywords{OCR \and modele wizyjno-językowe \and język polski \and dane syntetyczne \and znaki diakrytyczne}
\end{abstract}

%------------------------------------------------------------------------------
\section{Wprowadzenie}

\TODO{1 akapit: OCR, VLM, problem języków innych niż angielski.}

Yılmaz i~in.~\cite{yilmaz2026} zaproponowali syntetyczny benchmark dla języka tureckiego i~pokazali,
że znaki specyficzne dla języka są główną trudnością dla większości modeli. Język polski ma
dziewięć liter z~diakrytykami (ą, ć, ę, ł, ń, ó, ś, ź, ż), z~których część różni się od liter
łacińskich jedynie drobnym znakiem (ogonek, kreska, kropka).

\textbf{Cel.} Celem artykułu jest ocena, jak dokładnie dostępne lokalnie silniki OCR i~otwarty model
VLM rozpoznają tekst polski oraz jak na wynik wpływają polskie znaki, długość słowa, kontekst zdania,
typ tekstu i~zniekształcenia obrazu.

\textbf{Elementy nowości:}
(1) przeniesienie protokołu~\cite{yilmaz2026} na język polski wraz z~otwartym generatorem;
(2) uwzględnienie klasycznych silników z~modelem polskim (Tesseract \texttt{pol}, EasyOCR \texttt{pl}),
których autorzy nie badali;
(3) analiza błędów wyłącznie diakrytycznych i~ograniczeń słowników znaków;
(4) weryfikacja odtwarzalności badania bazowego.

%------------------------------------------------------------------------------
\section{Przegląd literatury}

\TODO{Krótko: Tesseract~\cite{smith2007}, PP-OCR~\cite{li2022ppocrv3}, TrOCR~\cite{li2023trocr},
VLM~\cite{bai2025qwen25vl}; dane syntetyczne~\cite{jaderberg2014,trdg}; benchmark turecki~\cite{yilmaz2026};
prace o~OCR dla polskiego (do uzupełnienia, max 10 pozycji łącznie).}

%------------------------------------------------------------------------------
\section{Materiał i metody}

\subsection{Zbiór danych}
Tekst źródłowy to \emph{Lalka, tom pierwszy}~\cite{wolnelektury} (domena publiczna). Pulę słów stanowią
wyrazy co najmniej trzyliterowe występujące w~tekście co najmniej dwa razy (8720 słów), a~pulę zdań
-- zdania liczące 4--8 słów, złożone wyłącznie z~liter i~interpunkcji (680 zdań).
Obrazy generujemy biblioteką TRDG 1.8.0~\cite{trdg} z~parametrami autorów badania bazowego: wysokość
64~px, marginesy 15~px, czarny tekst na białym tle. Tekst drukowany renderujemy 15 fontami (Open Sans,
Roboto, Source Sans 3), a~odręczny 19 popularnymi fontami odręcznymi z~Google Fonts. Każdy font
sprawdzamy w~tabeli \texttt{cmap} pod kątem wszystkich znaków korpusu, aby wykluczyć brakujące glify.

Dla każdego typu tekstu (drukowany, odręczny) tworzymy cztery kategorie:
(1) słowa z~polskimi znakami i~bez nich,
(2) słowa krótkie (${<}10$ znaków) i~długie (${>}10$),
(3) pojedyncze słowo i~zdanie,
(4) zniekształcenia: obraz czysty, rozmycie (promień 2), obrót (do $\pm16^\circ$), szum gaussowski
($\sigma=50$) i~niska rozdzielczość (wysokość 40~px), stosowane do tych samych słów.
Daje to 22 komórki po $n$ obrazów; \TODO{wyniki wstępne: $n=20$ (440 obrazów); docelowo $n=100$ (2200)}.
Losowanie jest deterministyczne (stałe ziarno), a~mniejsza próba jest podzbiorem większej.

\subsection{Modele}
\begin{itemize}
  \item \textbf{Tesseract} 5.5~\cite{smith2007}, model \texttt{pol}, tryb jednej linii (\texttt{--psm 7});
  \item \textbf{EasyOCR}~\cite{easyocr}, model \texttt{pl};
  \item \textbf{PaddleOCR} 2.9 (PP-OCRv4)~\cite{li2022ppocrv3}, model \texttt{pl} oraz \texttt{en}
    (ustawienie autorów badania bazowego);
  \item \textbf{TrOCR} large-printed~\cite{li2023trocr}, model angielski, jak w~badaniu bazowym,
    z~domyślnym limitem 21 tokenów wyjścia (co skraca dłuższe zdania, także u~autorów);
  \item \textbf{Qwen2.5-VL}~\cite{bai2025qwen25vl} \TODO{3B/7B}, prompt autorów z~,,Turkish''
    zamienionym na ,,Polish''.
\end{itemize}
Modele klasyczne uruchamiamy na CPU, \TODO{Qwen: CPU (3B) / Colab GPU (7B)}.

\subsection{Miary}
Niech $t_i$ oznacza tekst wzorcowy, a~$p_i$ predykcję modelu dla obrazu $i$, oba po normalizacji
(Unicode NFC, małe litery, scalone białe znaki). Główną miarą jest współczynnik błędów znakowych
\begin{equation}\label{eq:cer}
  \mathrm{CER}_i = \min\left(1,\ \frac{S_i + D_i + I_i}{N_i}\right),
\end{equation}
gdzie $S_i$, $D_i$, $I_i$ to liczby podstawień, usunięć i~wstawień znaków w~minimalnym ciągu edycji
(odległość Levenshteina~\cite{levenshtein1966}), a~$N_i = |t_i|$ to liczba znaków wzorca. Obcięcie do~1
odpowiada protokołowi~\cite{yilmaz2026}: predykcja dłuższa od wzorca nie może dać błędu powyżej 100\%.
Współczynnik błędów słów $\mathrm{WER}_i$ liczymy tym samym wzorem na ciągach słów i~podajemy tylko dla
zdań, bo dla pojedynczego słowa przyjmuje on wartości 0 lub~1.

Wynik dla zbioru obrazów $A$ (komórki, typu tekstu lub całego zbioru) to średnia
\begin{equation}\label{eq:mean}
  \overline{\mathrm{CER}}(A) = \frac{1}{|A|}\sum_{i\in A}\mathrm{CER}_i .
\end{equation}
Wpływ polskich znaków na błędy opisujemy odsetkiem \emph{błędów wyłącznie diakrytycznych}: wśród błędnych
predykcji $E=\{i: p_i\neq t_i\}$ liczymy te, które stają się poprawne po usunięciu diakrytyków
funkcją $s(\cdot)$ (ą$\to$a, ł$\to$l itd.):
\begin{equation}\label{eq:diacr}
  R_{\mathrm{diakr}} = \frac{|\{i\in E: s(p_i) = s(t_i)\}|}{|E|}.
\end{equation}

\subsection{Analiza statystyczna}
Niepewność średniej~\eqref{eq:mean} szacujemy jak w~\cite{yilmaz2026} nieparametrycznym bootstrapem:
z~wartości $\mathrm{CER}_i$ losujemy ze zwracaniem $B=1000$ prób o~liczności $|A|$, a~95\% przedział
ufności wyznaczają percentyle 2,5 i~97,5 rozkładu średnich z~prób. Wpływ czynnika (np.\ obecności
polskich znaków) mierzymy różnicą średnich między grupą z~czynnikiem $A$ i~grupą odniesienia $B$:
\begin{equation}\label{eq:delta}
  \Delta = \overline{\mathrm{CER}}(A) - \overline{\mathrm{CER}}(B),
\end{equation}
np.\ $A$ -- słowa z~polskimi znakami, $B$ -- słowa bez nich; dla zniekształceń $B$ to obrazy czyste.
Przedział ufności dla $\Delta$ wyznaczamy, losując niezależnie obie grupy; efekt uznajemy za istotny,
gdy przedział nie zawiera zera. Ziarno generatora losowego jest stałe, więc wyniki są powtarzalne.

%------------------------------------------------------------------------------
\section{Wyniki}

Tabela~\ref{tab:overall} przedstawia wyniki zbiorcze, a~tabela~\ref{tab:cells} -- wyniki w~komórkach.

\begin{table}
  \caption{Średni CER~\eqref{eq:mean} według typu tekstu z~95\% przedziałem ufności (bootstrap), WER dla
  zdań, odsetek błędów wyłącznie diakrytycznych~\eqref{eq:diacr} i~czas na obraz (CPU).}
  \label{tab:overall}
  \centering\small
  \begin{tabular}{lrrrcrrr}
\toprule
Model & \multicolumn{4}{c}{CER} & WER & Diakr. & s/obraz \\
 & druk & pismo & razem & 95\% CI & zdania & [\%] & \\
\midrule
Tesseract & .114 & .329 & .221 & [.191, .254] & .210 & 4 & 0.10 \\
EasyOCR & .085 & .257 & .171 & [.147, .199] & .392 & 8 & 0.12 \\
PaddleOCR (pl) & .111 & .261 & .186 & [.162, .213] & .489 & 24 & 0.05 \\
PaddleOCR (en) & .114 & .249 & .182 & [.157, .206] & .498 & 31 & 0.16 \\
TrOCR & .122 & .283 & .202 & [.182, .226] & .601 & 25 & 4.21 \\
\bottomrule
\end{tabular}

\end{table}

\begin{table}
  \caption{Średni CER w~komórkach: tekst drukowany / odręczny.}
  \label{tab:cells}
  \centering
  \resizebox{\textwidth}{!}{\begin{tabular}{lcccc}
\toprule
Podkategoria & Tesseract & EasyOCR & PaddleOCR & TrOCR \\
\midrule
z polskimi znakami & .006 / .149 & .006 / .192 & .180 / .383 & .204 / .268 \\
bez polskich znaków & .000 / .232 & .000 / .113 & .000 / .110 & .023 / .094 \\
\midrule
krótkie słowa & .000 / .110 & .000 / .084 & .106 / .243 & .131 / .216 \\
długie słowa & .000 / .131 & .004 / .078 & .070 / .087 & .070 / .117 \\
\midrule
słowo & .000 / .096 & .000 / .107 & .115 / .199 & .125 / .210 \\
zdanie & .004 / .196 & .025 / .250 & .080 / .107 & .187 / .249 \\
\midrule
czysty & .000 / .133 & .000 / .124 & .064 / .243 & .075 / .206 \\
rozmycie & .018 / .507 & .273 / .730 & .114 / .268 & .099 / .567 \\
obrót & .029 / .359 & .000 / .164 & .064 / .183 & .075 / .190 \\
szum & .723 / .954 & .000 / .164 & .069 / .216 & .075 / .229 \\
niska rozdzielczość & .471 / .752 & .627 / .819 & .363 / .837 & .281 / .761 \\
\bottomrule
\end{tabular}
}
\end{table}

\begin{table}
  \caption{Wpływ czynników: różnica średniego CER $\Delta$~\eqref{eq:delta} między grupą z~czynnikiem
  a~grupą odniesienia (dla zniekształceń: obrazy czyste). Pogrubienie: 95\% przedział ufności nie
  zawiera zera.}
  \label{tab:effects}
  \centering\small
  \begin{tabular}{lrrrr}
\toprule
Czynnik & Tesseract & EasyOCR & PaddleOCR & TrOCR \\
\midrule
polskie znaki & -.038 & +.043 & \textbf{+.226} & \textbf{+.178} \\
długie słowa & +.010 & -.001 & \textbf{-.096} & \textbf{-.080} \\
zdanie zamiast słowa & +.052 & \textbf{+.084} & \textbf{-.063} & +.051 \\
pismo odręczne & \textbf{+.215} & \textbf{+.172} & \textbf{+.150} & \textbf{+.160} \\
rozmycie & \textbf{+.196} & \textbf{+.439} & +.037 & \textbf{+.193} \\
obrót & \textbf{+.128} & +.020 & -.030 & -.008 \\
szum & \textbf{+.772} & +.020 & -.011 & +.012 \\
niska rozdzielczość & \textbf{+.545} & \textbf{+.661} & \textbf{+.446} & \textbf{+.381} \\
\bottomrule
\end{tabular}

\end{table}

\TODO{Opis wyników po kategoriach (tab.~\ref{tab:effects}): polskie znaki, długość, kontekst,
zniekształcenia, typ tekstu.}

%------------------------------------------------------------------------------
\section{Dyskusja}

\textbf{Słowniki znaków.} PaddleOCR dla języka polskiego korzysta ze wspólnego modelu łacińskiego,
którego słownik (\texttt{latin\_dict.txt}) zawiera litery ą, ć, ę, ł, ó, ale nie zawiera ń, ś, ź, ż.
Model nie może więc ich zwrócić i~zastępuje je literami podobnymi (np. ,,ludnošci'' zamiast
,,ludności''). Wyjaśnia to, dlaczego ustawienie \texttt{pl} daje wyniki zbliżone do \texttt{en}.

\textbf{Odtwarzalność badania bazowego.} Autorzy udostępnili dane, kod i~predykcje dla każdego obrazu.
Ponowne uruchomienie PaddleOCR na warstwowej próbie 330 obrazów dało predykcje identyczne
z~opublikowanymi w~99,7\% przypadków, tak samo jak TrOCR. Kod autorów zamiast wzoru~\eqref{eq:cer} liczy
\begin{equation}\label{eq:cer-authors}
  \mathrm{CER}^{\mathrm{kod}}_i = \frac{S_i + D_i + I_i}{|p_i|},
\end{equation}
czyli dzieli przez długość predykcji, a~nie wzorca (zamienione argumenty funkcji); tak policzone są
wartości w~opublikowanym pliku z~predykcjami. Liczby z~tekstu artykułu (np.\ TrOCR 0,38) odtwarza
natomiast inny wariant z~ich repozytorium: wzór~\eqref{eq:cer} po usunięciu spacji i~interpunkcji,
czego artykuł nie opisuje. Zmienia to ranking: GOT-OCR2, który gubi spacje, ma w~tym wariancie
CER 0,311, a~wzorem~\eqref{eq:cer} 0,391.

\TODO{Porównanie z~wynikami tureckimi (tylko jakościowo).}

%------------------------------------------------------------------------------
\section{Ograniczenia i dalsze prace}

Dane są syntetyczne, jednoliniowe i~na białym tle (tekstury papieru z~badania bazowego nie zostały
udostępnione); pominięto tekst w~scenie i~zdjęcia rzeczywiste. \TODO{Mała próba / jeden korpus /
mniejszy wariant Qwen.}

%------------------------------------------------------------------------------
\section{Wnioski}

\TODO{3--4 zdania.}

\subsubsection*{Dostępność danych i kodu.}
Generator, skrypty uruchomieniowe, wyniki dla każdego obrazu i~zbiór obrazów:
\TODO{adres repozytorium}. Kod na licencji MIT, dane na licencji CC BY 4.0.

\bibliographystyle{splncs04}
\bibliography{references}

\end{document}
